\chapter{Second-Order Linear Differential Equations}

In Chapter~37, we explored first-order differential equations: laws that tie
the rate of change of a quantity directly to the quantity itself, like cooling
coffee, radioactive decay, or a draining tank. A first-order equation takes your
current state and sets your velocity at that instant. But the physical world is governed by a deeper kind of law. Newton did
not write ``velocity is proportional to force.'' He wrote ``\emph{acceleration}
is proportional to force.'' Acceleration is the second derivative. The moment
you write down $F = ma$ for almost any real mechanical or electrical system -- a
mass on a spring, a pendulum, a swinging bridge, a tuned radio circuit, a
vibrating molecule -- you get a \emph{second-order} differential equation.

This chapter is about the most important and most tractable family of them: the
\emph{linear} ones with constant coefficients,
\[
  a\,y'' + b\,y' + c\,y \;=\; f(t).
\]
The entire theory hangs on one natural insight -- that $y = e^{rt}$
is the natural shape of a solution -- and this single guess, pushed all the way
through, unfolds into oscillation, damping, resonance, and the whole vocabulary
of vibration. Throughout, our running example is the \emph{harmonic oscillator},
a mass on a spring, because the harmonic oscillator provides the mathematical prototype
for second-order linear behavior.

\chapterline

\section{What a Second-Order Linear Equation Is}
\label{ch38:what}

A second-order ordinary differential equation (ODE) is an equation involving an
unknown function $y(t)$, its first derivative $y'$, and its second derivative
$y''$. In physical terms, it ties together an object's position, velocity, and
acceleration at each instant of time. It is \emph{linear} if $y$, $y'$, and
$y''$ each appear to the first power only, never multiplied together, never
inside a sine or a square root. The general second-order linear ODE with constant
coefficients is
\begin{equation}
  a\,y'' + b\,y' + c\,y = f(t),
  \qquad a \neq 0,
  \label{ch38:general}
\end{equation}
where $a$, $b$, $c$ are fixed numbers and $f(t)$ is a given function, the
\emph{forcing term} or \emph{input} -- an external push or driving signal
applied from the outside world. When $f(t) = 0$ the equation is called
\emph{homogeneous} (the system is left alone to move under its own internal
forces); otherwise it is \emph{inhomogeneous} (or \emph{forced}, driven by
outside intervention).

\begin{focusbox}
The plan of the chapter is direct: first solve the homogeneous
equation $a y'' + b y' + c y = 0$ completely (that is the ``free'' motion of the
system, what it does when nobody pushes it), then handle the forcing $f(t)$ by
adding one particular solution on top. The structure
\[
  y_{\text{general}} = \underbrace{y_h}_{\text{free motion}} + \underbrace{y_p}_{\text{response to forcing}}
\]
is the spine of everything that follows.
\end{focusbox}

\subsection*{Why second order means two constants}

A first-order equation comes with one arbitrary constant in its general
solution; you pin it down with one initial condition, usually $y(t_0)$. A
second-order equation comes with \emph{two} arbitrary constants, and you need
\emph{two} initial conditions to determine a specific solution. The natural pair
is
\[
  y(t_0) = y_0 \quad(\text{initial position}), \qquad
  y'(t_0) = v_0 \quad(\text{initial velocity}).
\]

\begin{intuitionbox}
This matches physical reality exactly. To predict where a thrown ball will be,
knowing the law of gravity is not enough -- you must know where you released it
\emph{and} how fast and in what direction. Position and velocity together form
the complete ``state'' of a Newtonian system. Two numbers in, a unique future
out. The two arbitrary constants of the general solution are precisely the slots
where that initial state gets plugged in.
\end{intuitionbox}

\begin{walkbox}
Before doing any algebra, pause and test the physical idea: why could a single
initial condition, such as knowing the pendulum passes through the bottom
$y(0) = 0$, never be enough to predict where it will be three seconds later?
Picture two grandfather clocks. In both clocks, the pendulum swings through the
bottom dead center ($y = 0$) at $t = 0$. But one pendulum was pulled far back and
is flying through the center at maximum speed, while the other was barely nudged
and is creeping past. Their positions match, but their velocities are completely
different -- and five seconds later, one is high on an upswing while the other has
already completed multiple cycles. Without velocity, the future remains
undetermined.
\end{walkbox}

\begin{historybox}
The leap from first to second order is the leap from Aristotle to Newton, but
the mathematical conquest of these equations belongs to Leonhard Euler.
Aristotle's physics was, in modern notation, \emph{first order}: he believed a
force produces a velocity ($b\,y' = f$), so remove the force and motion ceases.
Isaac Newton's revolution in the \emph{Principia} (1687) insisted that force
produces \emph{acceleration} ($m\,y'' = f$), giving birth to the second
derivative as the mathematical fingerprint of inertia -- the tendency of moving
matter to keep moving on its own.

Yet for fifty years after Newton, mathematicians were stuck. They had the
calculus of Newton and Leibniz, but every vibrating system -- a swinging
pendulum, an oscillating spring, a musical string -- was tackled through
individual geometric diagrams and ad-hoc approximations. Nobody possessed a
universal method to solve even the basic constant-coefficient equation
$a y'' + b y' + c y = 0$.

In 1739, in an exchange of letters with Johann Bernoulli, Euler found the key.
He proposed an \emph{ansatz} -- an educated guess for the algebraic shape of a
solution -- testing the exponential $y = e^{rt}$. Because differentiating an
exponential merely scales it by $r$, the differential equation collapsed into the
elementary quadratic $a r^2 + b r + c = 0$.

This stroke sparked controversy. When the discriminant $b^2 - 4ac$ was negative,
the roots $r = \alpha \pm i\beta$ involved $\sqrt{-1}$. Contemporaries
distrusted imaginary numbers as physical nonsense: how could an imaginary
exponent predict the real, tangible oscillation of a spring? Euler settled the
dispute by unveiling his formula $e^{i\theta} = \cos\theta + i\sin\theta$,
showing that complex exponentials trace rotations in the complex plane and combine
into ordinary real sines and cosines. Soon after, Jean le Rond d'Alembert resolved the
puzzle of repeated roots ($b^2 - 4ac = 0$), showing that the missing second
independent solution is $t e^{rt}$. A fierce debate followed between Euler,
d'Alembert, and Daniel Bernoulli over whether superpositions -- sums of pure sine
and cosine waves -- could describe an arbitrary plucked string with sharp corners,
a question only fully resolved a century later by Joseph Fourier.

Euler's exponential method unlocked acoustics, structural vibration in ships
and bridges, and later the entire theory of alternating-current (AC) circuits and
quantum wave mechanics. Yet the frontier did not close: equations with
time-dependent coefficients (like Bessel and Airy equations) rarely yield
elementary closed-form solutions, and nonlinear second-order equations (such as
the large-angle pendulum $y'' + \frac{g}{L}\sin y = 0$ and the gravitational
three-body problem) lead directly into deterministic chaos -- where exact formulas
fail and the long-term stability of orbits remains an active mathematical quest.
\end{historybox}

\section{The Harmonic Oscillator: Our Running Example}
\label{ch38:oscillator}

Picture a block of mass $m$ on a frictionless table, attached to a wall by a
spring (Figure~\ref{ch38:fig-spring}). Let $y(t)$ be the displacement of the
block from its rest position, positive to the right. Hooke's law says the spring
pulls back with a force proportional to how far you have stretched it, and
always toward rest: $F_{\text{spring}} = -k\,y$, where $k > 0$ is the spring
constant. Newton's second law $F = ma$ then reads
\[
  m\,y'' = -k\,y,
  \qquad\text{i.e.}\qquad
  m\,y'' + k\,y = 0.
\]

\begin{figure}[H]
\centering
\begin{tikzpicture}[scale=1.0]
  % wall
  \fill[pattern=north east lines] (-0.3,-0.7) rectangle (0,0.7);
  \draw[thick] (0,-0.7) -- (0,0.7);
  % spring (zigzag, stretched)
  \draw[thick,curveblue] (0,0) -- (0.3,0)
    -- (0.55,0.22) -- (0.95,-0.22) -- (1.35,0.22) -- (1.75,-0.22)
    -- (2.15,0.22) -- (2.55,-0.22) -- (2.8,0) -- (3.0,0);
  % mass
  \draw[thick,fill=blue!12] (3.0,-0.45) rectangle (4.2,0.45);
  \node at (3.6,0) {$m$};
  % ground
  \draw[thick] (-0.3,-0.7) -- (5.0,-0.7);
  % rest reference line
  \draw[dashed,gray] (2.4,-0.7) -- (2.4,1.8);
  \node[gray,above,font=\small] at (2.4,1.82) {rest ($y=0$)};
  % displacement arrow
  \draw[->,thick,tangentred] (2.4,1.35) -- (3.6,1.35) node[midway,above,font=\small]{displacement $y$};
  % restoring force arrow
  \draw[->,very thick,sourceblue] (3.6,0.68) -- (2.4,0.68) node[midway,above,font=\small]{$F = -ky$};
\end{tikzpicture}
\caption{A mass on a spring. When pulled to the right ($y > 0$), Hooke's law
generates a restoring force $F = -ky$ pulling back to the left toward rest. The
interplay of this restoring pull and the mass's inertia sustains the oscillation.}
\label{ch38:fig-spring}
\end{figure}

Divide by $m$ and write $\omega_0^2 = k/m$. The equation becomes the cleanest
second-order ODE in all of science:
\begin{equation}
  y'' + \omega_0^2\, y = 0.
  \label{ch38:sho}
\end{equation}
The constant $\omega_0 = \sqrt{k/m}$ is the \emph{natural frequency} -- the rate
in radians per second at which the spring wants to bounce when left entirely to
itself. Before we solve it formally, read the equation as a sentence: \emph{the
second derivative of $y$ is $-\omega_0^2$ times $y$ itself}. What function is
proportional to its own second derivative, with a \emph{minus} sign? Sine and
cosine. Differentiate $\cos(\omega_0 t)$ twice and you get
$-\omega_0^2 \cos(\omega_0 t)$. The form of the equation suggests its trigonometric solutions directly.

\begin{intuitionbox}
Why \emph{must} a spring oscillate, intuitively? When the block is displaced
right, the spring pulls it left. It accelerates left, races back through the rest
point with maximum speed, overshoots because of inertia, stretches the spring on
the left, gets pulled back right, and so on forever. The continuous interchange
between restoring force and inertia produces sinusoidal motion.
\end{intuitionbox}

If we add friction proportional to velocity (a dashpot, air resistance, the
block dragging on the table), the force law gains a damping resistance $-b\,y'$,
where $b > 0$ is the \emph{damping coefficient} measuring the strength of the
drag, and we get the \emph{damped} harmonic oscillator
\begin{equation}
  m\,y'' + b\,y' + k\,y = 0.
  \label{ch38:damped}
\end{equation}
If we also push the block with an external force $F(t)$ -- shaking the wall,
say -- we get the \emph{forced} oscillator $m y'' + b y' + k y = F(t)$. These
three equations, free, damped, and forced, will carry us through the entire
chapter. They are special cases of the general form~\eqref{ch38:general}, and
everything we prove about that form is best pictured as something happening to
this block.

\section{Superposition: Why Linearity Is Everything}
\label{ch38:superposition}

The word \emph{linear} is not decoration. It is the single property that makes
these equations solvable, and it deserves to be stated as a theorem.

\begin{theorem}[Superposition for the homogeneous equation]
\label{ch38:thm-superposition}
Suppose $y_1(t)$ and $y_2(t)$ are both solutions of the homogeneous equation
$a y'' + b y' + c y = 0$. Then for any constants $C_1, C_2$, the combination
$y = C_1 y_1 + C_2 y_2$ is also a solution.
\end{theorem}

\begin{proof}
Write $L[y] := a y'' + b y' + c y$ for the operation the equation performs on a
function. Differentiation is linear: the derivative of a sum is the sum of the
derivatives, and constants pull out. So
\[
  L[C_1 y_1 + C_2 y_2]
  = a(C_1 y_1 + C_2 y_2)'' + b(C_1 y_1 + C_2 y_2)' + c(C_1 y_1 + C_2 y_2).
\]
Distribute the derivatives and regroup by $y_1$ and $y_2$:
\[
  = C_1\underbrace{(a y_1'' + b y_1' + c y_1)}_{L[y_1] = 0}
  + C_2\underbrace{(a y_2'' + b y_2' + c y_2)}_{L[y_2] = 0}
  = C_1\cdot 0 + C_2 \cdot 0 = 0. \qedhere
\]
\end{proof}

\begin{happeningbox}
The operator $L[y] = a y'' + b y' + c y$ is a \emph{linear machine} -- a rule
that accepts an entire function as input and outputs a new function assembled
from its derivatives. Feed it a scaled sum, and out comes the scaled sum of the
outputs: $L[C_1 y_1 + C_2 y_2] = C_1 L[y_1] + C_2 L[y_2]$. This is the same
abstract property that makes matrices, integrals, and the derivative itself
behave so well. The set of all solutions of $L[y] = 0$ is therefore a
\emph{vector space} -- a collection of functions where adding any two solutions
or scaling by a number produces another valid solution. For a second-order
equation this space is two-dimensional, which is the deep reason there are
exactly two arbitrary constants. Finding the general solution means finding a
\emph{basis}: two independent solutions $y_1, y_2$ that act as coordinate axes,
spanning every possible free motion.
\end{happeningbox}

\begin{whybox}
Superposition is why our strategy works at all. Finding two linearly independent
solutions of the homogeneous equation is sufficient; linearity guarantees that
every solution is of the form $C_1 y_1 + C_2 y_2$. What ``genuinely different'' means
precisely -- linear independence -- is the subject we turn to after we have found
some solutions to be independent \emph{about}.
\end{whybox}

\section{Solving the Homogeneous Equation: Why \texorpdfstring{$e^{rt}$}{e\textasciicircum(rt)}?}
\label{ch38:characteristic}

We now want two solutions of
\begin{equation}
  a y'' + b y' + c y = 0.
  \label{ch38:homog}
\end{equation}
Where do we even start? The central idea of the chapter is worth a slow read.

\begin{happeningbox}
Look at what equation~\eqref{ch38:homog} asks for: a function whose second
derivative, first derivative, and self can be \emph{added up with constant weights
to cancel to zero}. For that cancellation to happen at every instant $t$, the
three functions $y''$, $y'$, $y$ had better all be the \emph{same shape} as each
other -- otherwise their sum could not vanish identically. We need a function
that reproduces itself under differentiation, up to a constant multiple. There is
essentially only one such function in all of mathematics:
\[
  \dv{}{t}\, e^{rt} = r\, e^{rt}.
\]
The exponential is the unique shape that differentiation cannot change, only
rescale. This is why $e^{rt}$ is directly suggested by the structure of the equation.
Every other method in this chapter is downstream of this one observation.
\end{happeningbox}

So we \emph{try} $y = e^{rt}$, with $r$ an unknown constant, and let the equation
tell us which $r$ works. Differentiating, $y' = r e^{rt}$ and $y'' = r^2 e^{rt}$.
Substitute into~\eqref{ch38:homog}:
\[
  a r^2 e^{rt} + b r e^{rt} + c e^{rt} = 0
  \quad\Longrightarrow\quad
  e^{rt}\,(a r^2 + b r + c) = 0.
\]
Now $e^{rt}$ is never zero, so we may divide it away. What survives is purely
algebraic -- the differential equation has collapsed into a quadratic:
\begin{equation}
  \boxed{\,a r^2 + b r + c = 0.\,}
  \label{ch38:characteristic-eq}
\end{equation}
This is the \emph{characteristic equation} (or \emph{auxiliary equation}) of the
ODE. We have traded a calculus problem for a high-school algebra problem. Its
roots $r$ are exactly the exponents that make $e^{rt}$ solve the differential
equation.

\begin{whybox}
Why is collapsing a differential equation into a polynomial legitimate -- and so
powerful? Because differentiating $e^{rt}$ multiplies by $r$, the operator
$a\,\frac{\dd^2}{\dd t^2} + b\,\frac{\dd}{\dd t} + c$ acts on $e^{rt}$ exactly the
way the polynomial $a r^2 + b r + c$ acts on the number $r$. Differentiation
``becomes'' multiplication by $r$. The hard analytic object (an operator on
functions) and the easy algebraic object (a quadratic in $r$) are two faces of
the same thing. This correspondence -- calculus turning into algebra on
exponentials -- is the seed of the Laplace transform and of huge parts of
engineering mathematics.
\end{whybox}

The quadratic~\eqref{ch38:characteristic-eq} has roots given by the familiar
formula
\[
  r = \frac{-b \pm \sqrt{b^2 - 4ac}}{2a}.
\]
Everything now hinges on the \emph{discriminant} $\Delta = b^2 - 4ac$ -- the
quantity inside the square root that decides whether the roots are two real
numbers, a single repeated real number, or a pair of complex numbers. Each case
yields a distinct physical behavior, and we examine them in turn.

\begin{walkbox}
Pause and test the intuition: if you face the equation $y'' - 9y = 0$, what two
exponential shapes does differentiation preserve?
Before writing down formulas, think: what number squared gives $+9$? It must be
$r = 3$ and $r = -3$. So the solutions are $e^{3t}$ and $e^{-3t}$. Notice the
stark contrast: with $y'' = +9y$, the curve accelerates away from zero instead of
oscillating back toward it. A mere plus sign in the equation turned an oscillator
into an exponential runaway.
\end{walkbox}

\section{Case 1: Two Distinct Real Roots \texorpdfstring{($\Delta > 0$)}{(Delta > 0)}}
\label{ch38:real-roots}

When $b^2 - 4ac > 0$, the characteristic equation has two different real roots
$r_1 \neq r_2$. Then $e^{r_1 t}$ and $e^{r_2 t}$ are both solutions, and -- as we
will confirm in \S\ref{ch38:wronskian} -- they are independent. Superposition
hands us the general solution
\begin{equation}
  y(t) = C_1 e^{r_1 t} + C_2 e^{r_2 t}.
  \label{ch38:real-general}
\end{equation}

Physically, this case is \emph{overdamping}: so much friction that the system
oozes back to rest without ever oscillating. Both roots are typically negative
(for a stable system), so both exponentials decay, just at different rates.

\begin{example}[Overdamped motion]
\label{ch38:ex-real}
Solve $y'' + 5y' + 6y = 0$ with $y(0) = 1$, $y'(0) = 0$.

\textbf{Step 1. Characteristic equation.} Replace $y'' \to r^2$, $y' \to r$,
$y \to 1$:
\[
  r^2 + 5r + 6 = 0 \;\Longrightarrow\; (r+2)(r+3) = 0
  \;\Longrightarrow\; r = -2,\ -3.
\]

\textbf{Step 2. General solution.} Two distinct real roots, so
\[
  y(t) = C_1 e^{-2t} + C_2 e^{-3t}.
\]

\textbf{Step 3. Apply initial conditions.} First, $y(0) = C_1 + C_2 = 1$.
Differentiate: $y'(t) = -2 C_1 e^{-2t} - 3 C_2 e^{-3t}$, so $y'(0) = -2C_1 - 3C_2
= 0$. Solve the pair:
\[
  C_1 + C_2 = 1, \qquad -2C_1 - 3C_2 = 0.
\]
From the first, $C_1 = 1 - C_2$; substitute: $-2(1-C_2) - 3C_2 = 0 \Rightarrow
-2 - C_2 = 0 \Rightarrow C_2 = -2$, hence $C_1 = 3$.

\textbf{Step 4. Answer.}
\[
  y(t) = 3 e^{-2t} - 2 e^{-3t}.
\]
At $t=0$ this is $3 - 2 = 1$ with zero slope, as required. As $t \to \infty$ both
terms vanish: the block creeps back to rest and never crosses zero. That is the
signature of overdamping.
\end{example}

\begin{figure}[H]
\centering
\begin{tikzpicture}
\begin{axis}[
  width=11cm, height=6cm,
  xlabel=$t$, ylabel=$y$,
  domain=0:4, samples=120,
  axis lines=left, xmin=0, xmax=4, ymin=-0.1, ymax=1.1,
  legend style={at={(0.5,1.03)}, anchor=south, legend columns=-1, font=\scriptsize},
  grid=both, grid style={gray, opacity=0.2},
]
  \addplot[curveblue, very thick] {3*exp(-2*x) - 2*exp(-3*x)};
  \addlegendentry{$y = 3e^{-2t}-2e^{-3t}$ (overdamped)}
  \addplot[gray, dashed, forget plot] {0};
\end{axis}
\end{tikzpicture}
\caption{The overdamped solution of Example~\ref{ch38:ex-real}. It returns
monotonically to rest -- too much friction to overshoot, so no oscillation.}
\label{ch38:fig-overdamped}
\end{figure}

\section{Case 3: Complex Roots \texorpdfstring{($\Delta < 0$)}{(Delta < 0)} -- Oscillation}
\label{ch38:complex-roots}

We jump to the complex case next, out of numerical order, because it is the heart
of the chapter: it is where oscillation lives. When $b^2 - 4ac < 0$, the square
root of a negative number forces the roots to be a \emph{complex-conjugate pair}
\[
  r = \alpha \pm i\beta,
  \qquad
  \alpha = \frac{-b}{2a}, \qquad
  \beta = \frac{\sqrt{4ac - b^2}}{2a} > 0.
\]
Formally, superposition still applies and the general solution is
$y = C_1 e^{(\alpha + i\beta)t} + C_2 e^{(\alpha - i\beta)t}$. But $y(t)$
describes the position of a real block; we want a \emph{real} answer, with no $i$
in sight. Euler's formula is what gets us there.

\begin{happeningbox}
Recall Euler's formula, $e^{i\theta} = \cos\theta + i\sin\theta$: as $\theta$
increases, $e^{i\theta}$ traces the unit circle in the complex plane, projecting
onto real cosines and sines. For a complex root $r = \alpha + i\beta$, the real
part $\alpha$ sets the growing or decaying amplitude envelope $e^{\alpha t}$,
while the imaginary part $\beta$ sets the oscillation frequency.
\end{happeningbox}

To extract the real solutions, factor out the common envelope:
\[
  e^{(\alpha \pm i\beta)t} = e^{\alpha t}\, e^{\pm i\beta t}
  = e^{\alpha t}\big(\cos\beta t \pm i \sin\beta t\big).
\]
Two particular real combinations come straight out by adding and subtracting.
Their average is the real part and (dividing by $2i$) their difference is the
imaginary part:
\[
  \tfrac12\big(e^{(\alpha+i\beta)t} + e^{(\alpha-i\beta)t}\big)
  = e^{\alpha t}\cos\beta t,
  \qquad
  \tfrac{1}{2i}\big(e^{(\alpha+i\beta)t} - e^{(\alpha-i\beta)t}\big)
  = e^{\alpha t}\sin\beta t.
\]
By superposition (Theorem~\ref{ch38:thm-superposition}), each of these is itself
a real solution. They are independent, so the real general solution is
\begin{equation}
  \boxed{\,y(t) = e^{\alpha t}\big(C_1 \cos\beta t + C_2 \sin\beta t\big).\,}
  \label{ch38:complex-general}
\end{equation}

\begin{whybox}
Why are we allowed to take real and imaginary parts and call them solutions?
Because the coefficients $a, b, c$ are real, the operator $L$ maps real functions
to real functions and complex to complex \emph{without mixing}: if $L[u + iv] = 0$
with $u, v$ real, then $L[u] + i\,L[v] = 0$, and a complex number is zero only
when both its real and imaginary parts vanish. So $L[u] = 0$ and $L[v] = 0$
separately. The real and imaginary parts of a complex solution are each, on their
own, real solutions.
\end{whybox}

\subsection*{The undamped oscillator, finally solved}

Return to the pure spring~\eqref{ch38:sho}, $y'' + \omega_0^2 y = 0$. The
characteristic equation is $r^2 + \omega_0^2 = 0$, giving $r = \pm i\omega_0$:
pure imaginary, $\alpha = 0$ and $\beta = \omega_0$. By~\eqref{ch38:complex-general},
\[
  y(t) = C_1 \cos\omega_0 t + C_2 \sin\omega_0 t.
\]
There is no envelope ($e^{0} = 1$), so the block oscillates forever with constant
amplitude -- exactly the perpetual duel between spring and inertia we pictured
earlier. The motion is often rewritten in \emph{amplitude-phase form}
\[
  y(t) = A\cos(\omega_0 t - \varphi),
  \qquad A = \sqrt{C_1^2 + C_2^2}, \quad \tan\varphi = \frac{C_2}{C_1},
\]
which makes the amplitude $A$ and the natural frequency $\omega_0$ stand out at a
glance.

\begin{example}[Undamped oscillation]
\label{ch38:ex-undamped}
Solve $y'' + 9y = 0$ with $y(0) = 2$, $y'(0) = -3$.

Characteristic equation $r^2 + 9 = 0 \Rightarrow r = \pm 3i$, so $\omega_0 = 3$
and
\[
  y(t) = C_1 \cos 3t + C_2 \sin 3t.
\]
Now $y(0) = C_1 = 2$. Differentiate: $y'(t) = -3C_1 \sin 3t + 3 C_2 \cos 3t$, so
$y'(0) = 3 C_2 = -3 \Rightarrow C_2 = -1$. Hence
\[
  y(t) = 2\cos 3t - \sin 3t.
\]
The amplitude is $A = \sqrt{2^2 + (-1)^2} = \sqrt 5 \approx 2.24$, and the motion
repeats with period $2\pi/3$ (Figure~\ref{ch38:fig-undamped}).
\end{example}

\begin{figure}[H]
\centering
\begin{tikzpicture}
\begin{axis}[
  width=12cm, height=5.5cm,
  xlabel=$t$, ylabel=$y$,
  domain=0:4.2, samples=200,
  axis lines=left, xmin=0, xmax=4.2, ymin=-2.6, ymax=2.6,
  legend style={at={(0.5,1.03)}, anchor=south, legend columns=-1, font=\scriptsize},
  grid=both, grid style={gray, opacity=0.2},
]
  \addplot[curveblue, very thick] {2*cos(deg(3*x)) - sin(deg(3*x))};
  \addlegendentry{$y = 2\cos 3t - \sin 3t$}
  \addplot[tangentred, dashed, thick, forget plot] {sqrt(5)};
  \addplot[tangentred, dashed, thick] {-sqrt(5)};
  \addlegendentry{amplitude $\pm\sqrt5$}
\end{axis}
\end{tikzpicture}
\caption{The undamped oscillation of Example~\ref{ch38:ex-undamped}: a pure sine
wave of constant amplitude $\sqrt5$ and period $2\pi/3$, riding forever between
the fixed lines $\pm\sqrt5$ because there is no friction to drain its energy.}
\label{ch38:fig-undamped}
\end{figure}

\subsection*{Underdamping: oscillation inside a shrinking envelope}

Now reinstate a \emph{little} friction in the damped oscillator~\eqref{ch38:damped},
$m y'' + b y' + k y = 0$, just enough that the discriminant stays negative. Then
$\alpha = -b/(2m) < 0$ and the solution is
\[
  y(t) = e^{-(b/2m)\,t}\big(C_1 \cos\beta t + C_2 \sin\beta t\big).
\]
The block still oscillates, but each swing is smaller than the last: the
oscillation rides inside a decaying envelope $\pm A e^{\alpha t}$. This is
\emph{underdamping} -- a plucked guitar string, a struck tuning fork, a car on
good shock absorbers nudged once. The frequency of the rings, $\beta$, is
slightly lower than the undamped $\omega_0$, because friction makes each cycle a
hair slower.

\begin{figure}[H]
\centering
\begin{tikzpicture}
\begin{axis}[
  width=12cm, height=6.5cm,
  xlabel=$t$, ylabel=$y$,
  domain=0:10, samples=200,
  axis lines=left, xmin=0, xmax=10, ymin=-1.1, ymax=1.1,
  legend style={at={(0.5,1.03)}, anchor=south, legend columns=-1, font=\scriptsize},
  grid=both, grid style={gray, opacity=0.2},
]
  \addplot[curveblue, very thick] {exp(-0.3*x)*cos(deg(2*x))};
  \addlegendentry{$y = e^{-0.3t}\cos 2t$ (underdamped)}
  \addplot[tangentred, dashed, thick, domain=0:10, forget plot] {exp(-0.3*x)};
  \addplot[tangentred, dashed, thick, domain=0:10] {-exp(-0.3*x)};
  \addlegendentry{envelope $\pm e^{-0.3t}$}
\end{axis}
\end{tikzpicture}
\caption{Underdamped motion: a sine wave squeezed inside the decaying envelope
$\pm e^{-0.3t}$. The system oscillates but loses amplitude every cycle. This is
the most common behavior of real vibrating systems.}
\label{ch38:fig-underdamped}
\end{figure}

\begin{aestheticbox}
Watch how three utterly different pictures -- a slow monotone crawl back to rest
(overdamping), an eternal pure sine wave (no damping), and a ringing wave fading
inside an envelope (underdamping) -- all fall out of \emph{one} quadratic by
merely changing the sign of $b^2 - 4ac$. Few places in mathematics show so clearly how a single
algebraic quantity governs qualitative behavior. The boundary case $\Delta = 0$
is the repeated-root case we treat next.
\end{aestheticbox}

\begin{figure}[H]
\centering
\begin{tikzpicture}[scale=0.82,>={Stealth[length=4pt]}]
  % Panel (a): Undamped (Center)
  \begin{scope}[shift={(0,0)}]
    \draw[fill=blue!3,draw=blue!15,rounded corners=4pt] (-2.5,-2.4) rectangle (2.5,2.75);
    \node[font=\small\bfseries,anchor=north] at (0,2.65) {(a) Undamped: Center};
    \draw[->,gray!50] (-2.1,0)--(2.1,0) node[right,font=\scriptsize] {$y$};
    \draw[->,gray!50] (0,-1.9)--(0,1.8) node[above,font=\scriptsize] {$v$};
    % Elliptical closed orbit (clockwise)
    \draw[curveblue,very thick] (0,0) ellipse (1.6cm and 1.2cm);
    \draw[->,curveblue,ultra thick] (0,1.2) -- (0.3,1.2);
    \draw[->,curveblue,ultra thick] (0,-1.2) -- (-0.3,-1.2);
    \node[font=\scriptsize,align=center] at (0,-2.1) {Closed orbit: energy conserved};
  \end{scope}

  % Panel (b): Underdamped (Stable Spiral)
  \begin{scope}[shift={(5.7,0)}]
    \draw[fill=green!3,draw=green!15,rounded corners=4pt] (-2.5,-2.4) rectangle (2.5,2.75);
    \node[font=\small\bfseries,anchor=north] at (0,2.65) {(b) Underdamped: Spiral};
    \draw[->,gray!50] (-2.1,0)--(2.1,0) node[right,font=\scriptsize] {$y$};
    \draw[->,gray!50] (0,-1.9)--(0,1.8) node[above,font=\scriptsize] {$v$};
    % Inward spiral (clockwise)
    \draw[targetgreen,very thick,samples=150,domain=0:18.84,variable=\th]
      plot ({1.7*exp(-0.16*\th)*cos(deg(\th))}, {-1.3*exp(-0.16*\th)*sin(deg(\th))});
    \filldraw[targetgreen] (0,0) circle (2pt);
    \draw[->,targetgreen,ultra thick] ({1.7*exp(-0.16*1.57)*cos(deg(1.57))}, {-1.3*exp(-0.16*1.57)*sin(deg(1.57))}) -- ++(0.3,0);
    \node[font=\scriptsize,align=center] at (0,-2.1) {Inward spiral: energy drained};
  \end{scope}

  % Panel (c): Overdamped (Stable Node)
  \begin{scope}[shift={(11.4,0)}]
    \draw[fill=orange!3,draw=orange!15,rounded corners=4pt] (-2.5,-2.4) rectangle (2.5,2.75);
    \node[font=\small\bfseries,anchor=north] at (0,2.65) {(c) Overdamped: Node};
    \draw[->,gray!50] (-2.1,0)--(2.1,0) node[right,font=\scriptsize] {$y$};
    \draw[->,gray!50] (0,-1.9)--(0,1.8) node[above,font=\scriptsize] {$v$};
    % Eigenvector directions
    \draw[dashed,gray!60,thick] (-1.8,0.9) -- (1.8,-0.9) node[pos=0.08,left,font=\scriptsize,gray!80!black] {fast};
    \draw[dashed,gray!60,thick] (-1.8,-1.2) -- (1.8,1.2) node[pos=0.92,right,font=\scriptsize,gray!80!black] {slow};
    % Trajectories approaching along slow axis
    \draw[->,quotientorange,very thick] (-1.6,1.4) .. controls (-0.8,0.2) and (-0.5,-0.2) .. (-0.05,-0.03);
    \draw[->,quotientorange,very thick] (1.6,-1.4) .. controls (0.8,-0.2) and (0.5,0.2) .. (0.05,0.03);
    \filldraw[quotientorange] (0,0) circle (2pt);
    \node[font=\scriptsize,align=center] at (0,-2.1) {Monotonic creep: no rotation};
  \end{scope}
\end{tikzpicture}
\caption{The three vibrational regimes viewed in displacement--velocity phase space $(y, v)$ where $v=y'$.
(a) Undamped motion is a closed periodic orbit (Center) cycling forever between potential and kinetic energy.
(b) Underdamped motion is an inward spiral (Spiral Sink), losing amplitude on each cycle as friction extracts energy.
(c) Overdamped motion is a direct approach (Nodal Sink), collapsing along the slow eigenspace into rest without a single oscillation.}
\label{ch38:fig-phase-space-comparison}
\end{figure}

\section{Case 2: A Repeated Real Root \texorpdfstring{($\Delta = 0$)}{(Delta = 0)}}
\label{ch38:repeated-root}

The borderline case $b^2 - 4ac = 0$ is the strangest and most instructive. The
quadratic has a single (double) root
\[
  r = -\frac{b}{2a},
\]
and the recipe $e^{rt}$ gives us only \emph{one} solution, $y_1 = e^{rt}$. But a
second-order equation needs \emph{two} independent solutions to span its
two-dimensional solution space. We are one solution short. Where does the second
one come from?

\begin{trapbox}
A tempting wrong move is to write $y = C_1 e^{rt} + C_2 e^{rt} = (C_1 + C_2)
e^{rt}$. That is \emph{not} a two-parameter family at all -- it collapses to a
single constant times $e^{rt}$. You have one solution masquerading as two. The
two pieces are not independent, so the initial-value problem has no hope of
matching arbitrary $y(0)$ \emph{and} $y'(0)$. A genuinely new solution is
required.
\end{trapbox}

\begin{happeningbox}
Here is the real reason the missing solution is $t\,e^{rt}$. Think of the
repeated root as the limit of two distinct roots $r$ and $r + \eps$ that are
sliding together as $\eps \to 0$. For distinct roots we have the two solutions
$e^{rt}$ and $e^{(r+\eps)t}$. Their difference, divided by $\eps$, is also a
solution by superposition:
\[
  \frac{e^{(r+\eps)t} - e^{rt}}{\eps}.
\]
But this is exactly a difference quotient -- in the limit $\eps \to 0$ it becomes
the \emph{derivative of $e^{rt}$ with respect to $r$}:
\[
  \pdv{}{r}\, e^{rt} = t\,e^{rt}.
\]
The second solution is the rate at which $e^{rt}$ changes as you nudge the
exponent. When two exponential solutions collide, the slot the second one
vacated is filled by ``$t$ times the first.''
\end{happeningbox}

Let us verify directly that $y_2 = t\,e^{rt}$ works when $r = -b/(2a)$ is a double
root, so that $a r^2 + b r + c = 0$ \emph{and} (since $r$ is the vertex of the
parabola) $2 a r + b = 0$. Compute the derivatives of $y_2 = t e^{rt}$:
\[
  y_2' = e^{rt} + r t\, e^{rt} = (1 + rt)e^{rt},
  \qquad
  y_2'' = r e^{rt} + r(1+rt)e^{rt} = (2r + r^2 t)e^{rt}.
\]
Substitute into $a y_2'' + b y_2' + c y_2$ and collect by powers of $t$:
\[
  e^{rt}\Big[\underbrace{(a r^2 + b r + c)}_{=\,0}\, t
  + \underbrace{(2 a r + b)}_{=\,0}\Big] = 0.
\]
Both bracketed quantities vanish precisely because $r$ is a \emph{repeated} root.
So $t e^{rt}$ is a bona fide second solution. The general solution in the
repeated-root case is therefore
\begin{equation}
  \boxed{\,y(t) = (C_1 + C_2 t)\,e^{rt}.\,}
  \label{ch38:repeated-general}
\end{equation}

Physically this is \emph{critical damping}: the exact knife's-edge amount of
friction that brings the system back to rest as fast as possible \emph{without}
overshooting. A screen door closer, a sensitive galvanometer, and a well-tuned
car suspension are all designed to sit at or near this critical value -- any less
and they would bounce, any more and they would respond sluggishly.

\begin{example}[Critical damping]
\label{ch38:ex-repeated}
Solve $y'' + 4y' + 4y = 0$ with $y(0) = 1$, $y'(0) = 1$.

Characteristic equation $r^2 + 4r + 4 = (r+2)^2 = 0$, a double root $r = -2$. So
\[
  y(t) = (C_1 + C_2 t)e^{-2t}.
\]
Then $y(0) = C_1 = 1$. Differentiate with the product rule:
\[
  y'(t) = C_2 e^{-2t} - 2(C_1 + C_2 t)e^{-2t} = \big(C_2 - 2C_1 - 2 C_2 t\big)e^{-2t}.
\]
So $y'(0) = C_2 - 2C_1 = 1$, giving $C_2 = 1 + 2 = 3$. The solution is
\[
  y(t) = (1 + 3t)e^{-2t}.
\]
Notice it rises briefly (the $3t$ wins early) before the exponential drags it
down to rest -- a single gentle bump, no oscillation. That is critical damping in
one picture.
\end{example}

\begin{figure}[H]
\centering
\begin{tikzpicture}
\begin{axis}[
  width=12cm, height=6cm,
  xlabel=$t$, ylabel=$y$,
  domain=0:5, samples=150,
  axis lines=left, xmin=0, xmax=5, ymin=-0.1, ymax=1.5,
  legend style={at={(0.97,0.97)},anchor=north east}, grid=both, grid style={gray!15},
]
  \addplot[structurepurple, very thick] {(1+3*x)*exp(-2*x)};
  \addlegendentry{critical: $(1+3t)e^{-2t}$}
  \addplot[gray, dashed, forget plot] {0};
\end{axis}
\end{tikzpicture}
\caption{Critically damped motion from Example~\ref{ch38:ex-repeated}: one bump,
then a clean return to rest with no overshoot. This is the fastest non-oscillating
decay possible.}
\label{ch38:fig-critical}
\end{figure}

\section{When Are Two Solutions Really Independent? The Wronskian}
\label{ch38:wronskian}

We have repeatedly claimed that pairs like $\{e^{r_1 t}, e^{r_2 t}\}$ or
$\{e^{rt}, t e^{rt}\}$ are ``independent,'' so that $C_1 y_1 + C_2 y_2$ captures
\emph{all} solutions. Let us pin down what ``independent'' means, since that is
exactly what guarantees we have not missed any solutions.

Two solutions $y_1, y_2$ are \emph{linearly independent} if neither is a constant
multiple of the other -- equivalently, if $C_1 y_1 + C_2 y_2 = 0$ for all $t$
forces $C_1 = C_2 = 0$. The clean test is the \emph{Wronskian} determinant
\[
  W(y_1, y_2)(t) = \det\begin{pmatrix} y_1 & y_2 \\ y_1' & y_2' \end{pmatrix}
  = y_1 y_2' - y_2 y_1'.
\]

\begin{whybox}
Why does this determinant detect independence? Suppose we want constants with
$C_1 y_1 + C_2 y_2 = 0$ for all $t$. Differentiating gives a second equation
$C_1 y_1' + C_2 y_2' = 0$. Together they are a $2\times 2$ linear system for
$(C_1, C_2)$:
\[
  \begin{pmatrix} y_1 & y_2 \\ y_1' & y_2' \end{pmatrix}
  \begin{pmatrix} C_1 \\ C_2 \end{pmatrix} = \begin{pmatrix} 0 \\ 0 \end{pmatrix}.
\]
If the determinant $W$ is nonzero at even one point, the only solution is $C_1 =
C_2 = 0$ -- so $y_1, y_2$ are independent. The Wronskian is precisely the
determinant that decides whether the ``cancel to zero'' system has a nontrivial
solution. A nonzero Wronskian is the certificate that our two solutions genuinely
span a two-dimensional space, with no overlap.
\end{whybox}

\begin{example}[The exponential pair is independent]
\label{ch38:ex-wronskian}
For $y_1 = e^{r_1 t}$, $y_2 = e^{r_2 t}$ with $r_1 \neq r_2$,
\[
  W = e^{r_1 t}\cdot r_2 e^{r_2 t} - e^{r_2 t}\cdot r_1 e^{r_1 t}
    = (r_2 - r_1)\,e^{(r_1 + r_2)t}.
\]
Since $r_1 \neq r_2$ and the exponential never vanishes, $W \neq 0$ everywhere.
The two solutions are independent, confirming~\eqref{ch38:real-general}. For the
repeated-root pair $\{e^{rt}, t e^{rt}\}$ a similar computation gives $W = e^{2rt}
\neq 0$, so that pair is independent too.
\end{example}

\begin{formalbox}
The payoff is the \emph{Existence and Uniqueness Theorem} for constant-coefficient
linear ODEs: given any initial data $y(t_0) = y_0$, $y'(t_0) = v_0$, there is one
and only one solution. Because the Wronskian of our basis $\{y_1, y_2\}$ is
nonzero, the system
\[
  C_1 y_1(t_0) + C_2 y_2(t_0) = y_0, \qquad
  C_1 y_1'(t_0) + C_2 y_2'(t_0) = v_0
\]
always has a unique solution $(C_1, C_2)$. So every initial-value problem is
solved exactly once by our two-parameter family -- nothing is missed, nothing is
double-counted. This is the rigorous backbone under all the case work above.
\end{formalbox}

\section{Forcing the System: Inhomogeneous Equations}
\label{ch38:forcing}

Now switch the forcing back on:
\begin{equation}
  a y'' + b y' + c y = f(t).
  \label{ch38:inhomog}
\end{equation}
Physically, $f(t)$ is an external push: a periodic shove on the spring, a voltage
applied to a circuit, gravity loading a beam. The structure theorem that organizes
everything is this.

\begin{theorem}[Structure of the general solution]
\label{ch38:thm-structure}
Let $y_p$ be \emph{any one} particular solution of the inhomogeneous
equation~\eqref{ch38:inhomog}, and let $y_h = C_1 y_1 + C_2 y_2$ be the general
solution of the associated homogeneous equation $a y'' + b y' + c y = 0$. Then
the general solution of~\eqref{ch38:inhomog} is
\[
  y = y_h + y_p = C_1 y_1 + C_2 y_2 + y_p.
\]
\end{theorem}

\begin{proof}
First, $y_h + y_p$ solves~\eqref{ch38:inhomog}: applying the linear operator $L$,
$L[y_h + y_p] = L[y_h] + L[y_p] = 0 + f = f$. Conversely, let $Y$ be \emph{any}
solution. Then $L[Y - y_p] = L[Y] - L[y_p] = f - f = 0$, so $Y - y_p$ solves the
homogeneous equation and must equal some $C_1 y_1 + C_2 y_2$. Hence $Y = y_h +
y_p$. Every solution has this form, and every function of this form is a solution.
\end{proof}

\begin{intuitionbox}
Read the structure physically. The homogeneous part $y_h$ is the system's own
free vibration -- whatever it would do if you set it going and walked away. The
particular part $y_p$ is its forced response -- the steady motion the external
push drives it into. The true motion is the \emph{sum}: the driven response, plus
whatever leftover free vibration the initial conditions demand. In a damped
system $y_h$ decays away (it is called the \emph{transient}), and after a while
only the forced $y_p$ survives (the \emph{steady state}). This is why, long after
you start shaking a damped spring, it settles into the rhythm of your hand and
forgets how it started.
\end{intuitionbox}

\section{Undetermined Coefficients: Guessing \texorpdfstring{$y_p$}{y\_p} Intelligently}
\label{ch38:undetermined}

To finish the inhomogeneous problem we need just one particular solution $y_p$.
For the common forcing functions -- polynomials, exponentials, sines and
cosines -- there is a beautifully direct method: \emph{guess a $y_p$ of the same
shape as $f$, with unknown coefficients, and let the equation solve for them.}

\begin{whybox}
Why should a same-shape guess work? Because differentiation keeps these special
functions inside their own families. Differentiate a polynomial and you get a
polynomial; differentiate $e^{kt}$ and you get a multiple of $e^{kt}$;
differentiate $\cos\omega t$ and you stay within $\{\cos\omega t,
\sin\omega t\}$. Since $L[y]$ only differentiates and scales, if $y_p$ is built
from such a family, then $L[y_p]$ lands in the same family. To match a forcing
$f$ from that family, we only need to tune the coefficients. The method is not
luck; it exploits the fact that these forcings live in finite-dimensional spaces
that $L$ maps into themselves.
\end{whybox}

The table of first guesses:
\begin{center}
\begin{tabular}{ll}
\toprule
Forcing $f(t)$ & Trial $y_p$ \\
\midrule
constant $k$ & $A$ \\
polynomial degree $n$ & $A_n t^n + \dots + A_1 t + A_0$ \\
$e^{kt}$ & $A e^{kt}$ \\
$\cos\omega t$ or $\sin\omega t$ & $A\cos\omega t + B\sin\omega t$ \\
$e^{kt}\cos\omega t$ & $e^{kt}(A\cos\omega t + B\sin\omega t)$ \\
\bottomrule
\end{tabular}
\end{center}
Always include \emph{both} a sine and a cosine when the forcing has either,
because differentiation turns one into the other.

\begin{example}[Polynomial forcing]
\label{ch38:ex-undetermined-poly}
Solve $y'' - y' - 2y = 4t^2$.

\textbf{Homogeneous part.} $r^2 - r - 2 = (r-2)(r+1) = 0$, so $r = 2, -1$ and
$y_h = C_1 e^{2t} + C_2 e^{-t}$.

\textbf{Guess.} The forcing is a degree-2 polynomial, so try $y_p = A t^2 + B t +
C$. Then $y_p' = 2At + B$ and $y_p'' = 2A$. Substitute:
\[
  \underbrace{2A}_{y_p''} - \underbrace{(2At + B)}_{y_p'} - 2\underbrace{(At^2 + Bt + C)}_{y_p}
  = 4t^2.
\]
Collect by powers of $t$:
\[
  -2A\,t^2 + (-2A - 2B)\,t + (2A - B - 2C) = 4t^2 + 0t + 0.
\]

\textbf{Match coefficients.}
\[
  t^2:\ -2A = 4 \Rightarrow A = -2; \qquad
  t^1:\ -2A - 2B = 0 \Rightarrow B = 2; \qquad
  t^0:\ 2A - B - 2C = 0 \Rightarrow C = -3.
\]
So $y_p = -2t^2 + 2t - 3$, and the general solution is
\[
  y = C_1 e^{2t} + C_2 e^{-t} - 2t^2 + 2t - 3.
\]
\end{example}

\begin{example}[Sinusoidal forcing -- a driven oscillator]
\label{ch38:ex-undetermined-sin}
Solve $y'' + 4y = 5\cos t$ (an undamped spring with $\omega_0 = 2$, pushed at
frequency $1$).

\textbf{Homogeneous part.} $r^2 + 4 = 0 \Rightarrow r = \pm 2i$, so $y_h = C_1
\cos 2t + C_2 \sin 2t$.

\textbf{Guess.} Forcing $5\cos t$; try $y_p = A\cos t + B\sin t$. Then $y_p'' =
-A\cos t - B\sin t$, and
\[
  y_p'' + 4 y_p = (-A + 4A)\cos t + (-B + 4B)\sin t = 3A\cos t + 3B\sin t.
\]
Match to $5\cos t + 0\sin t$: $3A = 5 \Rightarrow A = 5/3$, and $3B = 0 \Rightarrow
B = 0$. So $y_p = \tfrac53 \cos t$, and
\[
  y = C_1 \cos 2t + C_2 \sin 2t + \tfrac{5}{3}\cos t.
\]
The driving frequency $1$ differs from the natural frequency $2$, so the response
stays bounded -- the system simply hums along at the driving frequency with a
modest amplitude. Watch what happens in the next section when those two
frequencies coincide.
\end{example}

\subsection*{The resonance snag, and how to fix the guess}

There is one trap, and it is the most physically important phenomenon in the
chapter. \emph{What if the trial $y_p$ is already a solution of the homogeneous
equation?} Then substituting it gives $L[y_p] = 0 \neq f$, and the guess cannot
possibly match a nonzero forcing.

\begin{trapbox}
If your trial solution duplicates a homogeneous solution, do \emph{not} just give
up -- and do not try strange new shapes. The fix mirrors the repeated-root cure:
multiply the trial by $t$ (or by $t^2$ if the homogeneous root is doubled). For
forcing $e^{kt}$ when $k$ is already a characteristic root, use $y_p = A t e^{kt}$.
For forcing $\cos\omega_0 t$ at the natural frequency, use $y_p = t(A\cos\omega_0 t
+ B\sin\omega_0 t)$. The extra factor of $t$ plays the same role as in the
repeated-root case: it generates an independent solution when the trial form is
already an element of the homogeneous solution space.
\end{trapbox}

\section{Resonance: When Forcing Hits the Natural Frequency}
\label{ch38:resonance}

Take the undamped oscillator and push it at \emph{exactly} its own natural
frequency:
\begin{equation}
  y'' + \omega_0^2 y = F_0 \cos\omega_0 t.
  \label{ch38:resonance-eq}
\end{equation}
The naive guess $A\cos\omega_0 t + B\sin\omega_0 t$ is a homogeneous solution, so
by the rule above we multiply by $t$ and try
\[
  y_p = t\,(A\cos\omega_0 t + B\sin\omega_0 t).
\]
Carrying out the differentiation (product rule twice) and substituting, the
algebra delivers
\[
  y_p = \frac{F_0}{2\omega_0}\, t \sin\omega_0 t.
\]
The defining feature is the factor of $t$ \emph{outside} the sine. The amplitude
of the oscillation is $\dfrac{F_0}{2\omega_0}\,t$ -- it grows without bound,
linearly in time. This is \emph{resonance}.

\begin{happeningbox}
Why does pushing at the natural frequency make the amplitude explode? Timing.
When you drive at $\omega_0$, every push is perfectly synchronized with the
oscillator's own motion -- you always shove in the direction it is already
moving, like pumping a swing exactly in rhythm. Each cycle you add a little
energy and \emph{none} is wasted fighting the motion. Energy accumulates cycle
after cycle, so the amplitude climbs and climbs. The factor of $t$ in $t\sin
\omega_0 t$ is the mathematical signature of energy piling up with no escape. Off
resonance, your pushes sometimes help and sometimes oppose, the books roughly
balance, and the amplitude stays bounded -- as we saw in
Example~\ref{ch38:ex-undetermined-sin}.
\end{happeningbox}

\begin{figure}[H]
\centering
\begin{tikzpicture}
\begin{axis}[
  width=12cm, height=6.5cm,
  xlabel=$t$, ylabel=$y$,
  xlabel style={at={(1,0)}, anchor=south west},
  ylabel style={anchor=south west},
  domain=0:20, samples=200,
  axis lines=left, xmin=0, xmax=20, ymin=-11, ymax=11.5,
  legend style={at={(0.05, 0.95)}, anchor=north west}, grid=both, grid style={gray, opacity=0.2},
]
  \addplot[curveblue, very thick] {0.5*x*sin(deg(x))};
  \addlegendentry{$y = 0.5\, t\sin t$ (resonant response)}
  \addplot[tangentred, dashed, thick, forget plot] {0.5*x};
  \addplot[tangentred, dashed, thick] {-0.5*x};
  \addlegendentry{envelope $\pm 0.5 t$}
\end{axis}
\end{tikzpicture}
\caption{Resonance: driving at the natural frequency produces oscillations whose
amplitude grows linearly in $t$, bounded only by the widening envelope $\pm
\tfrac12 t$. No matter how small the push, given enough time the response becomes
arbitrarily large.}
\label{ch38:fig-resonance}
\end{figure}

\begin{aestheticbox}
Resonance is not a curiosity; it is a force of engineering life and death. A file
of soldiers is told to break step when crossing a bridge, lest their cadence
happen to match the bridge's natural frequency and pump it to destruction. The
Tacoma Narrows Bridge tearing itself apart in 1940 is the textbook (if
technically more subtle) cautionary image. On the constructive side, resonance is
how a radio works: a tunable circuit is built to resonate at one broadcast
frequency, amplifying that station hugely while ignoring all the others. A
microwave oven drives water molecules near their resonance; an MRI machine
resonates atomic nuclei. That factor of $t$ marks the exact frequency at which
energy builds up cycle after cycle.
\end{aestheticbox}

\begin{applicationbox}
Real systems always have a little damping, so the amplitude does not literally go
to infinity -- the decaying transient and the energy lost to friction eventually
balance the energy you pump in, and the response levels off at a large but finite
amplitude. The closer the driving frequency is to $\omega_0$, and the smaller the
damping $b$, the taller that resonant peak. Engineers plot the steady-state
amplitude against driving frequency and get a sharp spike at $\omega_0$ -- the
\emph{resonance curve}. Designing a structure means either keeping its natural
frequencies far from any expected driving frequency, or adding enough damping to
tame the peak. The undamped, infinitely growing solution above is the idealized
limit that warns you where the danger lies.
\end{applicationbox}

\begin{example}[Pure resonance, worked]
\label{ch38:ex-resonance}
Solve $y'' + y = 2\cos t$ with $y(0) = 0$, $y'(0) = 0$.

Here $\omega_0 = 1$ and the forcing frequency is also $1$: resonance. The
homogeneous solution is $y_h = C_1 \cos t + C_2 \sin t$. Since $\cos t$ solves the
homogeneous equation, multiply the trial by $t$:
\[
  y_p = t(A\cos t + B\sin t).
\]
Differentiate twice:
\[
  y_p' = A\cos t + B\sin t + t(-A\sin t + B\cos t),
\]
\[
  y_p'' = 2(-A\sin t + B\cos t) + t(-A\cos t - B\sin t).
\]
Then $y_p'' + y_p = 2(-A\sin t + B\cos t)$ (the $t$-terms cancel because $\cos t,
\sin t$ solve the homogeneous equation). Match to $2\cos t$: $2B = 2 \Rightarrow B
= 1$ and $-2A = 0 \Rightarrow A = 0$. So $y_p = t\sin t$, and
\[
  y = C_1\cos t + C_2 \sin t + t\sin t.
\]
Apply $y(0) = 0 \Rightarrow C_1 = 0$. Then $y' = C_2\cos t + \sin t + t\cos t$, so
$y'(0) = C_2 = 0$. The clean answer is
\[
  y(t) = t\sin t,
\]
an oscillation whose swings grow without bound -- the picture of
Figure~\ref{ch38:fig-resonance} with envelope $\pm t$.
\end{example}

\section{A Unified Worked Lab}
\label{ch38:lab}

\begin{labbox}
Take the damped, forced oscillator
\[
  y'' + 2y' + 5y = 10\cos t, \qquad y(0) = 0,\ y'(0) = 0,
\]
and carry it through end to end. The skeleton:
\begin{enumerate}
  \item \textbf{Homogeneous solution.} Characteristic equation $r^2 + 2r + 5 =
  0$. Discriminant $4 - 20 = -16 < 0$, so $r = -1 \pm 2i$. Underdamped:
  \[
    y_h = e^{-t}(C_1 \cos 2t + C_2 \sin 2t).
  \]
  \item \textbf{Particular solution.} Forcing $10\cos t$; the frequency $1$ is
  not a homogeneous frequency, so no resonance, no extra $t$. Try $y_p = A\cos t +
  B\sin t$ and show
  \[
    y_p'' + 2y_p' + 5y_p = (4A + 2B)\cos t + (4B - 2A)\sin t.
  \]
  Match to $10\cos t + 0\sin t$: solve $4A + 2B = 10$, $-2A + 4B = 0$. Confirm
  $A = 2$, $B = 1$, so $y_p = 2\cos t + \sin t$.
  \item \textbf{Combine.} $y = e^{-t}(C_1\cos 2t + C_2\sin 2t) + 2\cos t + \sin t$.
  \item \textbf{Initial conditions.} Use $y(0)=0$ to get $C_1 + 2 = 0$, so $C_1 =
  -2$. Now differentiate the full $y$ and set $y'(0) = 0$. Since
  \[
    y'(0) = -C_1 + 2C_2 + 1 = 0,
  \]
  substituting $C_1 = -2$ gives $2 + 2C_2 + 1 = 0$, so $C_2 = -\tfrac32$. Final
  answer:
  \[
    y(t) = -e^{-t}\big(2\cos 2t + \tfrac32 \sin 2t\big) + 2\cos t + \sin t.
  \]
  \item \textbf{Interpret.} Identify the transient $-e^{-t}\big(2\cos 2t +
  \tfrac32\sin 2t\big)$ (gone within a few seconds) and the steady state $2\cos t
  + \sin t$ (the lasting response,
  oscillating at the \emph{driving} frequency $1$, not the damped frequency
  $2$). Sketch both and watch the transient fade.
\end{enumerate}
Do every step with pencil before reading on; the differentiation in step~2 is the
only fiddly part, and getting it by hand cements the method.
\end{labbox}

\begin{figure}[H]
\centering
\begin{tikzpicture}
\begin{axis}[
  width=12cm, height=6.2cm,
  xlabel=$t$, ylabel=$y$,
  xlabel style={at={(1,0)}, anchor=south west},
  ylabel style={anchor=south west},
  domain=0:10, samples=200,
  axis lines=left, xmin=0, xmax=10, ymin=-3.2, ymax=4.8,
  legend style={at={(0.98,0.98)},anchor=north east,font=\footnotesize},
  grid=both, grid style={gray, opacity=0.2}
]
  \addplot[curveblue, very thick] {-exp(-x)*(2*cos(deg(2*x)) + 1.5*sin(deg(2*x))) + 2*cos(deg(x)) + sin(deg(x))};
  \addlegendentry{full solution $y = y_h + y_p$}
  \addplot[quotientorange, dashed, thick] {2*cos(deg(x)) + sin(deg(x))};
  \addlegendentry{steady state $y_p = 2\cos t + \sin t$}
  \addplot[targetgreen, dotted, very thick] {-exp(-x)*(2*cos(deg(2*x)) + 1.5*sin(deg(2*x)))};
  \addlegendentry{transient $y_h = -e^{-t}(2\cos 2t + 1.5\sin 2t)$}
  \addplot[gray, dashed, forget plot] {0};
\end{axis}
\end{tikzpicture}
\caption{Decomposition of the worked lab solution into its transient and steady-state components. The transient free response $y_h(t)$ (dotted green) decays exponentially within a few cycles as damping drains initial energy. The steady-state forced response $y_p(t)$ (dashed orange) oscillates indefinitely at the driving frequency. The actual motion $y(t) = y_h + y_p$ (solid blue) starts from rest at the origin ($y(0)=0, y'(0)=0$) and rapidly synchronizes with the driving force.}
\label{ch38:fig-lab-decomposition}
\end{figure}

\begin{connectionbox}
Everything here generalizes upward and sideways. \emph{Upward in order:} an
$n$th-order constant-coefficient linear ODE has a degree-$n$ characteristic
polynomial, and its $n$ roots (real, repeated, or complex) build the solution by
exactly the rules of this chapter -- $e^{rt}$ for each root, an extra $t$ for each
repetition, $e^{\alpha t}\cos\beta t$ pairs for complex roots. \emph{Sideways into
systems:} a single second-order equation is equivalent to a $2\times 2$
\emph{system} of first-order equations (set $v = y'$), and there the
characteristic equation reappears as the eigenvalue equation of a matrix. The
roots $r$ become eigenvalues; oscillation becomes complex eigenvalues; damping
becomes negative real parts.

In engineering applications, structures rarely consist of an isolated
oscillator. When multiple masses or electrical loops interact, coupled physical
components require simultaneous differential equations. The displacement of each
mass exerts a force on the next, meaning no single equation can be solved in
isolation. In Chapter~39, we step into that world, discovering how matrix
algebra and eigenvectors unravel systems of simultaneous ODEs into independent
modes of vibration.
\end{connectionbox}

\section{Summary}
\label{ch38:summary}

\begin{focusbox}
The whole method, on one card.

\textbf{To solve} $a y'' + b y' + c y = f(t)$:
\begin{enumerate}
  \item Solve the homogeneous equation. Write the characteristic equation $a r^2
  + b r + c = 0$ and find its roots.
  \begin{itemize}
    \item Distinct real $r_1, r_2$: \quad $y_h = C_1 e^{r_1 t} + C_2 e^{r_2 t}$
    (overdamped).
    \item Repeated real $r$: \quad $y_h = (C_1 + C_2 t)e^{rt}$ (critically damped).
    \item Complex $\alpha \pm i\beta$: \quad $y_h = e^{\alpha t}(C_1\cos\beta t +
    C_2\sin\beta t)$ (oscillation; underdamped if $\alpha < 0$).
  \end{itemize}
  \item If $f \neq 0$, find one particular $y_p$ by undetermined coefficients --
  guess the shape of $f$, multiplying by $t$ if the guess collides with a
  homogeneous solution (resonance).
  \item General solution $y = y_h + y_p$. Apply the two initial conditions
  $y(t_0), y'(t_0)$ last, to fix $C_1, C_2$.
\end{enumerate}
The central idea throughout: $e^{rt}$ converts the linear differential equation into an
algebraic quadratic equation, whose discriminant classifies the unforced motion as
overdamped, critically damped, or underdamped, while forcing at the natural frequency produces resonance.
\end{focusbox}

In Chapter~39, we link several variables together and study \textbf{systems of
differential equations}, where matrix eigenvalues play the exact role that
characteristic roots played here.

\section{Exercises}
\label{ch38:exercises}

\begin{exercise}
\label{ch38:ex-1}
Classify each by its discriminant as overdamped, critically damped, or
underdamped, then write the general solution.
(a) $y'' + 7y' + 12y = 0$. \quad
(b) $y'' + 6y' + 9y = 0$. \quad
(c) $y'' + 2y' + 10y = 0$.
\end{exercise}

\begin{exercise}
\label{ch38:ex-2}
Solve the initial-value problem $y'' + \omega_0^2 y = 0$, $y(0) = y_0$, $y'(0) =
v_0$, for general $\omega_0, y_0, v_0$. Then show the amplitude of the resulting
oscillation is $A = \sqrt{y_0^2 + v_0^2/\omega_0^2}$. Interpret each term
physically.
\end{exercise}

\begin{exercise}
\label{ch38:ex-3}
For the equation $y'' + 4y' + 4y = 0$ (a double root at $r = -2$), verify by
direct substitution that $y_2 = t e^{-2t}$ is a solution, and compute the
Wronskian $W(e^{-2t}, t e^{-2t})$ to confirm the two solutions are independent.
\end{exercise}

\begin{exercise}
\label{ch38:ex-4}
Solve $y'' - 3y' + 2y = 0$ with $y(0) = 0$, $y'(0) = 1$. Then describe, in one
sentence each, how the solution behaves as $t \to \infty$ and as $t \to -\infty$.
\end{exercise}

\begin{exercise}
\label{ch38:ex-5}
Use undetermined coefficients to find a particular solution of $y'' - y = 3e^{2t}$.
Why does the plain guess $A e^{2t}$ work here without an extra factor of $t$?
\end{exercise}

\begin{exercise}
\label{ch38:ex-6}
Find a particular solution of $y'' + y = 4\sin 2t$. The driving frequency is $2$
and the natural frequency is $1$ -- explain why no resonance occurs and no extra
$t$ is needed, then solve.
\end{exercise}

\begin{exercise}
\label{ch38:ex-7}
The forcing in $y'' - y = 4e^{t}$ matches a homogeneous solution ($r = 1$ is a
root of $r^2 - 1 = 0$). Use the corrected trial $y_p = A t e^{t}$ to find a
particular solution. This is the exponential analogue of resonance.
\end{exercise}

\begin{exercise}
\label{ch38:ex-8}
Pure resonance. Solve $y'' + 9y = 6\cos 3t$ with $y(0) = 0$, $y'(0) = 0$. Identify
the growing-amplitude term and state the envelope within which the solution
oscillates.
\end{exercise}

\begin{exercise}
\label{ch38:ex-9}
Damped driven oscillator. For $y'' + 2y' + 2y = 10\cos t$, find the steady-state
(particular) solution and its amplitude. Which part of the full solution is the
transient, and roughly how long until it is negligible?
\end{exercise}

\begin{exercise}
\label{ch38:ex-10}
A spring-mass system has $m = 1$, spring constant $k = 4$, and damping $b$. (a)
For which value of $b$ is the system critically damped? (b) Sketch (qualitatively)
the response to the initial data $y(0) = 1$, $y'(0) = 0$ for $b$ slightly below,
equal to, and slightly above that critical value, and explain the difference in
words.
\end{exercise}

\begin{exercise}
\label{ch38:ex-11}
Reduction to a system. Given $y'' + b y' + c y = 0$, set $v = y'$ and write the
equivalent pair of first-order equations for $(y, v)$. Show that looking for
solutions of the form $(y, v) = \vec{u}\,e^{rt}$ recovers exactly the
characteristic equation $r^2 + b r + c = 0$. (This is the bridge to the next
chapter.)
\end{exercise}

\begin{exercise}
\label{ch38:ex-12}
Conceptual. Explain, without solving any equation, why a second-order linear ODE
needs exactly two initial conditions while a first-order one needs only one.
Connect your answer to the dimension of the homogeneous solution space and to the
physical idea of a system's ``state.''
\end{exercise}
